% !TEX root = ../main.tex
\chapter{Data Science and AI for Biomedical Portfolios}\label{ch:biomed}
\begin{paracol}{2}

\begin{sources}
\begin{tabularx}{\linewidth}{@{}l l L@{}}
\toprule
\textbf{Lecture} & \textbf{Speaker} & \textbf{Content}\\
\midrule
2024 L18 & A.~W.~Lo (MIT Sloan; guest) & Biomedicine at an inflection point (gene therapy); the ``valley of death''; risk versus uncertainty; Sharpe ratio and the Madoff example; one anti-cancer drug as an investment; a portfolio of 150 independent projects; rare diseases, correlation and BridgeBio; a fusion ``mega fund''; questions (patents, expertise, COVID-19 vaccines).\\
\bottomrule
\end{tabularx}

\smallskip
\textbf{Files.} Transcript only: \file{Lecture 18 - Applying Data Science and Artificial Intelligence to Managing Biomedical Portfolios.txt}. The same text exists as
\file{18642-lecture-18-version-5_transcript.pdf} in \file{18.642-fall-2024/static_resources}. There are no slides, no
problem set and no other file for this lecture; everything said about the slides below is therefore unknown, and the
chapter is built on the spoken text.
\textbf{Problem sets.} None. The exercises at the end are all \emph{not from the course}.
\textbf{Not from the lecture.} Every derivation marked ``Reconstruction'' (correlated projects, binomial and
beta-binomial counts, securitisation, ML for trial success) is my own mathematics behind the numbers quoted by
Lo; the numerics were checked by a short Python script. The 2024 L1 slides list a past final-project paper,
\emph{A Megafund Approach to Pediatric Oncology Drug Development}, which uses exactly this framework.
\end{sources}

\section*{Motivation}
Prof.\ Lo is a financial economist (director of the MIT Laboratory for Financial Engineering, principal investigator
at CSAIL) \ts{24}{18}{0:00}\ts{24}{18}{2:07}. His message is not about biology: it is that a \emph{financing structure} can change the
risk--reward profile of an investment, and that the required mathematics is a few lines of probability. A single
drug-development project is a lottery ticket that nobody wants: costly, slow, and very likely to fail. A
portfolio of many such tickets whose outcomes are only weakly correlated has a much better ratio of expected
return to risk, and that is what makes it fundable. The chapter follows the lecture and then adds the
mathematics that the lecture only quotes:
\begin{itemize}
  \item the Sharpe ratio as the quantity to be raised, and the two ways to raise it (\cref{ch12:sec-sharpe});
  \item the single-project calculation with its numbers, and the $\sqrt N$ law for independent projects
        (\cref{ch12:sec-single,ch12:sec-portfolio});
  \item \emph{what breaks when projects are correlated}: the variance of a sum, the diversification limit,
        counts of successes, and parameter uncertainty as a hidden common factor
        (\cref{ch12:sec-corr,ch12:sec-counts}); this is the point the lecturer stresses with rare diseases and
        with the fusion question;
  \item a sketch of securitisation (debt and equity tranches) and of predicting trial success from data
        (\cref{ch12:sec-tranche,ch12:sec-ml});
  \item the second example of the lecture, fusion energy, and a numerical erratum (\cref{ch12:sec-fusion}).
\end{itemize}
Portfolio variance is developed in general in \cref{ch:portfolio}; the covariance formula for sums is in
\cref{ch:probability}.

% ==================================================================
\section{The setting: a valley of death}\label{ch12:sec-setting}

The lecture opens with two clinical stories, both single-gene diseases treated by a one-time gene therapy
(a correct copy of the gene is carried by a virus into the cells): Leber's congenital amaurosis, a blindness
caused by one DNA ``typo'' \ts{24}{18}{4:07}, and Canavan disease, in which the myelin insulation of nerve
cells degrades and children rarely live beyond about ten years \ts{24}{18}{6:13}\ts{24}{18}{9:19}. Lo uses them, together with
the ``convergence'' of life sciences, physical sciences and engineering and the various ``omics''
\ts{24}{18}{13:36}, to argue that science is advancing fast. The bottleneck he identifies is \emph{economics}
\ts{24}{18}{14:39}.

The specific obstacle is the \emph{valley of death}: the gap between academic experiments and the first
human clinical trials, where money is hard to raise \ts{24}{18}{15:42}. He explains it with the distinction, familiar in economics, between
\begin{itemize}
  \item \emph{risk}: randomness whose probability law can be parametrised (normal, gamma, $\chi^2$), and
  \item \emph{uncertainty}: ``unknown unknowns'', with no data from which to assess probabilities
        \ts{24}{18}{16:47}.
\end{itemize}
Drug development has three features that investors must digest: it costs hundreds of millions to billions
of dollars per approved drug, the probability of success is low, and the process takes 10 to 15 years
\ts{24}{18}{17:50}. Compared with writing an app, all three are extreme.

\begin{intuition}[Risk versus uncertainty, in physicist's terms]
Risk is a stochastic model with known parameters (a Hamiltonian with a heat bath at known temperature);
uncertainty is not knowing the Hamiltonian. \cref{ch12:sec-counts} shows that the two are not unrelated:
uncertainty about a parameter that all projects share acts exactly like a common random field coupling them,
and it limits diversification just as correlation does.
\end{intuition}

% ==================================================================
\section{Sharpe ratio: what investors want}\label{ch12:sec-sharpe}

Lo's classroom exercise shows four cumulative-return curves, unnamed, from 2008 onward, and asks which one
to choose \ts{24}{18}{18:54}\ts{24}{18}{19:59}. The four assets are revealed as Treasury bills
(safe, about zero real return), the S\&P~500, Pfizer, and the Fairfield Sentry fund, the feeder of Madoff's
Ponzi scheme, which was chosen by most because it looked like the best trade-off \ts{24}{18}{20:04}\ts{24}{18}{21:04}. The moral
is that investors want a high expected return per unit of risk, formalised by the Sharpe ratio.

\begin{definition}[Sharpe ratio]\label{ch12:def-sharpe}
For an investment with return $R$, risk-free return $r_f$ over the same horizon, the \emph{Sharpe ratio} is
\[
  \mathrm{SR}=\frac{\E[R]-r_f}{\sqrt{\Var R}}.
\]
\end{definition}
The lecture quotes $\mathrm{SR}\approx0.33$ for the S\&P~500 and $0.34$ for Pfizer, and an ``order of magnitude
higher'' figure, on paper, for the Ponzi scheme \ts{24}{18}{22:07}. Lo adds that the Sharpe ratio of
biomedical investing has been declining for decades, not because the numerator is unattractive but
because risk and uncertainty grow faster as the science becomes more sophisticated \ts{24}{18}{23:07}.

There are, jokes Lo, two ways to increase a ratio: raise the numerator or lower the denominator
(or both) \ts{24}{18}{24:11}. Better science and engineering act on both. What is under-used in biomedicine is
\emph{new financing and business models}, and this is where mathematicians can contribute: constructing
structures whose risk--reward trade-off differs from that of the underlying assets
\ts{24}{18}{25:14}.

\begin{finance}[Why the Madoff curve fooled people]
A smooth, upward curve has a tiny sample standard deviation, so its Sharpe ratio is enormous; a Sharpe ratio
that looks too good is itself a warning sign. In the lecture the point is behavioural (``we all want
high-yielding, low-risk assets'') and it sets up the target: not a lottery ticket, but a pooled structure with a
respectable Sharpe ratio \ts{24}{18}{22:07}.
\end{finance}

% ==================================================================
\section{One anti-cancer drug as an investment}\label{ch12:sec-single}

Lo poses the question as a fund. Raise \$200 million, wait 10 years, and with probability $95\%$ return nothing
\ts{24}{18}{25:14}. The back-of-the-envelope numbers for an oncology project are \ts{24}{18}{26:19}\ts{24}{18}{27:22}:
\begin{itemize}
  \item cost of the clinical trials: about \$200 million, up front;
  \item 10 years to complete them;
  \item probability of success: a bit under $5\%$ historically ($3.4\%$ in oncology), rounded to $p=5\%$;
  \item if successful, revenues of \$2 billion per year in years 11--20 (the patent lifetime is 20 years, so the
        exclusivity runs until then), equivalent to a single payment of \$12.3 billion in year~10.
\end{itemize}
The equivalence is a present-value computation (\cref{sec:discounting}). The lecture does not state the discount
rate, but the numbers are reproduced by $r=10\%$:
\begin{equation}\label{ch12:eq-annuity}
  2\times\sum_{k=1}^{10}\frac{1}{1.1^{\,k}}=2\times\frac{1-1.1^{-10}}{0.1}=2\times6.1446=12.29\approx12.3
  \quad\text{(billion, at year 10).}
\end{equation}
(This is a \emph{Reconstruction}: only the \$12.3 billion is quoted.)

\begin{example}[Reconstruction: mean and volatility of the single project]\label{ch12:ex-single}
Let $G$ be the gross ten-year return, $G=12.3/0.2=61.5$ with probability $p=0.05$ and $0$ otherwise. Then
\[
  \E G=pG_0=3.075,\qquad \Var G=G_0^2\,p(1-p),\qquad \sqrt{\Var G}=61.5\sqrt{0.0475}=13.40,
\]
with $G_0=61.5$. The annualised expected return is $3.075^{1/10}-1\approx11.9\%$, the ``about $12\%$'' of the lecture. The
annualised standard deviation, obtained by scaling the ten-year standard deviation by $\sqrt{10}$, is
$13.40/\sqrt{10}=4.24$, that is $424\%$, the ``$423.5\%$'' of the lecture \ts{24}{18}{28:23} (the small
difference comes from rounding of the inputs). The Sharpe ratio with $r_f\approx0$ is $0.12/4.24\approx0.03$,
``close to zero'': hardly anyone would invest.
\end{example}
Lo's conclusion is that few investors take this bet \ts{24}{18}{28:23}. The scaling $\sqrt{10}$ used for the
volatility is the usual rule for independent yearly increments; how the expected return is annualised is a convention
(the simple average $207.5\%/10\approx21\%$ would give another number), so the pair $(12\%,423.5\%)$ is best taken as
the lecturer's stated input to the next step.

% ==================================================================
\section{A portfolio of independent projects}\label{ch12:sec-portfolio}

Instead of one project, invest in $N=150$ at once: this needs $150\times\$200\text{M}=\$30$ billion, a sum that
Lo defers (``assume we have it'') \ts{24}{18}{29:26}. If the projects are independent and identically distributed
the expected return of the portfolio equals that of each project, while
the standard deviation is divided by $\sqrt N$: $423.5\%/\sqrt{150}=34.6\%$ \ts{24}{18}{30:30}. The Sharpe ratio is
then about $12\%/34.6\%\approx0.35$, ten times the single-project value.

\begin{proposition}[Reconstruction: equal-weight portfolio of $N$ i.i.d.\ projects]\label{ch12:prop-iid}
Let $R_1,\dots,R_N$ be i.i.d.\ with mean $\mu$ and variance $\sigma^2$ (returns per dollar invested) and let
$R_P=\frac1N\sum_i R_i$. Then $\E R_P=\mu$, $\Var R_P=\sigma^2/N$ and $\mathrm{SR}_N=\sqrt N\,\mathrm{SR}_1$.
\end{proposition}
\begin{proof}
Linearity of $\E$ gives the mean. For the variance, independence gives
$\Var\bigl(\sum_iR_i\bigr)=\sum_i\Var R_i=N\sigma^2$, and dividing the sum by $N$ divides the variance by $N^2$.
\end{proof}
With the numbers of \cref{ch12:ex-single}: $\mathrm{SR}_1=0.0283$, $\mathrm{SR}_{150}=0.0283\sqrt{150}=0.347$.
Lo phrases the step as ``the variance of the sum is the sum of the variances'' and ``the risk goes down by the square
root of the number of projects''; the second phrase refers to the standard deviation \ts{24}{18}{29:26}\ts{24}{18}{30:30}.

\begin{intuition}[Central-limit bookkeeping]
Ten thousand independent spins have relative fluctuations of order $10^{-2}$ of their mean: this is the same
$1/\sqrt N$ law that makes pressure well defined. The economics of ``a \$30 billion fund'' is the statement that
the fund is a large-$N$ system whose fluctuations are small, provided the couplings between the components are
small; the rest of the chapter is about that proviso.
\end{intuition}

Lo then asks where \$30 billion would come from, and answers that if the fund is attractive enough, small
investors can supply it in aggregate (part of a retirement account or of a summer-internship salary)
\ts{24}{18}{30:30}. This is the second half of the strategy: the same mathematics that makes the fund attractive
also justifies broad retail participation, and \cref{ch12:sec-tranche} sketches how such capital can be
structured.

% ==================================================================
\section{Correlation: the real constraint}\label{ch12:sec-corr}

The i.i.d.\ assumption is flagged by Lo himself (``I'll come back to that assumption later'')
\ts{24}{18}{29:26}. He returns to it with \emph{rare diseases} \ts{24}{18}{31:30}. There are roughly $12{,}000$
of them, each with a small patient population, but more than 30 million Americans are affected in
aggregate. Development is comparatively cheap because trials can be small (he mentions a
Leber's-amaurosis approval on 30--40 patients) and the target is a known gene \ts{24}{18}{32:31}\ts{24}{18}{33:36}. Above all
the diseases are so different that the successes and failures of different programmes are barely
correlated. To see why this matters, recall the variance of a sum.

\begin{proposition}[Variance of a portfolio of $N$ correlated projects]\label{ch12:prop-corr}
Let $R_1,\dots,R_N$ have common variance $\sigma^2$ and common pairwise correlation $\rho$. Then
\begin{equation}\label{ch12:eq-varcorr}
  \Var R_P=\frac{1}{N^2}\Bigl(\sum_i\Var R_i+\sum_{i\neq j}\Cov(R_i,R_j)\Bigr)
  =\sigma^2\Bigl(\frac1N+\frac{N-1}{N}\,\rho\Bigr),
\end{equation}
so that
\begin{equation}\label{ch12:eq-sharpecorr}
  \mathrm{SR}_N=\frac{\mathrm{SR}_1}{\sqrt{\,1/N+(1-1/N)\rho\,}}\ \xrightarrow[N\to\infty]{}\ \frac{\mathrm{SR}_1}{\sqrt\rho}.
\end{equation}
\end{proposition}
\begin{proof}
The double sum $\sum_{i\neq j}$ has $N(N-1)$ ordered pairs, each equal to $\rho\sigma^2$. Hence
$\Var R_P=\frac{1}{N^2}\bigl(N\sigma^2+N(N-1)\rho\sigma^2\bigr)$, which is~\eqref{ch12:eq-varcorr}.
The Sharpe ratio uses the same expected return, whence~\eqref{ch12:eq-sharpecorr}.
\end{proof}
This is the statement of the lecture: with $N$ projects there are $N(N-1)/2$ pairwise covariances, ``a lot of them'',
and in a typical investment portfolio they are almost all positive \ts{24}{18}{33:36}\ts{24}{18}{34:38}. Only when they vanish
does the risk fall as $1/\sqrt N$. The relevant limit of \eqref{ch12:eq-varcorr} is the
\emph{undiversifiable} floor $\sigma^2\rho$: no number of projects removes the common component.

\begin{table}[ht]
\centering\small
\caption{Reconstruction: standard deviation and Sharpe ratio of an equal-weight portfolio for the single-project
inputs of \cref{ch12:ex-single} (mean $12\%$, standard deviation $423.5\%$, $r_f=0$), computed from
\eqref{ch12:eq-varcorr}.}\label{ch12:tab-corr}
\begin{tabular}{@{}lrrrr@{}}
\toprule
& \multicolumn{2}{c}{$N=15$} & \multicolumn{2}{c}{$N=150$}\\
\cmidrule(lr){2-3}\cmidrule(l){4-5}
$\rho$ & std.\ dev. & SR & std.\ dev. & SR\\
\midrule
$0$    & $109\%$ & $0.11$ & $34.6\%$ & $0.35$\\
$0.01$ & $117\%$ & $0.10$ & $54.6\%$ & $0.22$\\
$0.02$ & $124\%$ & $0.10$ & $69.0\%$ & $0.17$\\
$0.05$ & $143\%$ & $0.08$ & $100\%$ & $0.12$\\
$0.10$ & $169\%$ & $0.07$ & $138\%$ & $0.09$\\
\bottomrule
\end{tabular}
\end{table}

\begin{figure}[ht]
\centering
\begin{tikzpicture}
\begin{axis}[width=0.82\linewidth,height=6.2cm,xmode=log,xlabel={number of projects $N$},
  ylabel={Sharpe ratio $\mathrm{SR}_N$},xmin=1,xmax=1000,ymin=0,ymax=1.0,
  legend pos=north west,legend cell align=left,legend style={font=\footnotesize},grid=both,
  grid style={black!10},samples=50,domain=1:1000,every axis plot/.append style={thick}]
\addplot[defcol]{0.0283*sqrt(x)};
\addlegendentry{$\rho=0$}
\addplot[thmcol]{0.0283/sqrt(1/x+(1-1/x)*0.02)};
\addlegendentry{$\rho=0.02$}
\addplot[fincol]{0.0283/sqrt(1/x+(1-1/x)*0.05)};
\addlegendentry{$\rho=0.05$}
\addplot[intcol]{0.0283/sqrt(1/x+(1-1/x)*0.2)};
\addlegendentry{$\rho=0.2$}
\end{axis}
\end{tikzpicture}
\caption{Reconstruction: Sharpe ratio of an equal-weight portfolio of $N$ projects with pairwise correlation $\rho$,
from \eqref{ch12:eq-sharpecorr} with $\mathrm{SR}_1=0.0283$. For $\rho=0$ it grows as $\sqrt N$; for $\rho>0$ it saturates
at $\mathrm{SR}_1/\sqrt\rho$.}\label{ch12:fig-sharpe}
\end{figure}

\begin{coursenote}[How many projects, how much money?]
Lo's claim is that because rare-disease programmes are nearly uncorrelated, one needs ``10 or 20''
projects rather than 150, hence \$400--500 million rather than \$30 billion \ts{24}{18}{34:38}. That is not obtained by scaling the
\$200 million per project of the oncology example: 15 such projects would cost \$3 billion. It also assumes that rare-disease
projects are cheaper and have a different profile than the oncology numbers used above, which the lecture does
not quantify. The reported input is a Monte Carlo simulation with industry numbers published in his paper,
giving about $22\%$ per year and a Sharpe ratio above~1 \ts{24}{18}{35:38}; the parameters are not in the
lecture, so \cref{ch12:tab-corr} should \emph{not} be read as a reproduction of that result: with the oncology
inputs and $N=15$ a Sharpe ratio of $0.11$ is all one gets, even with independence.
\end{coursenote}

\begin{finance}[The BridgeBio story]\label{ch12:fin-bridge}
A former student of Lo (a biotech venture capitalist) reran the simulations with his own parameters,
then with further modifications; they ran about 35 versions over six months, after which he quit his job to
found a company \ts{24}{18}{36:43}. BridgeBio Pharma came out of stealth in 2017 with a private-equity investor
(KKR) taking part in a biotech for the first time; it raised about \$40 million after a friends-and-family round,
then \$135 million, then \$300 million (about \$400--500 million in total), had about 15 projects, went public and raised a further
\$750 million, and traded at about \$4.6 billion of market capitalisation on the day of the lecture
\ts{24}{18}{37:45}\ts{24}{18}{38:48}. By his account about 20 projects were in the clinic, two drugs had been approved in 2021
and the Canavan gene therapy shown earlier belongs to the company \ts{24}{18}{39:52}. Lo's remark, worth
keeping: finance cannot turn a bad drug into a good one, but the wrong financing can destroy a good one; science
comes first \ts{24}{18}{41:57}. (Facts about the company are as stated in class; they were not checked.)
\end{finance}

% ==================================================================
\section{Counting successes}\label{ch12:sec-counts}

For a project financed at year 0 the natural random variable is a Bernoulli success indicator
$X_i\in\{0,1\}$ with $\Prob(X_i=1)=p$. This section is a \emph{Reconstruction} (not from the course) of what
the number of successes looks like when the indicators are independent or correlated.

\paragraph{Independent projects.}
$K=\sum_iX_i\sim\mathrm{Bin}(N,p)$. The probability of at least one success is $1-(1-p)^N$: with $p=5\%$ it is
$54\%$ for $N=15$ and $99.95\%$ for $N=150$ (the chance of nothing is $0.95^{150}=4.6\times10^{-4}$). If every success returns
$G_0=61.5$ times a single project's cost and every project has the same cost, the portfolio (total cost $N$ units) returns
$G_0K$ units, so the portfolio is at least at break-even when $K\ge N/G_0$. For $N=150$ this means $K\ge3$
($150/61.5=2.44$), which has probability $98.2\%$; for $N=15$ it means $K\ge1$, probability $54\%$.

\paragraph{Correlated outcomes.}
Two Bernoulli variables with the same $p$ have $\Cov(X_i,X_j)=\rho\,p(1-p)$, and then
\begin{equation}\label{ch12:eq-varK}
  \Var K=Np(1-p)\bigl(1+(N-1)\rho\bigr),
\end{equation}
by the same counting argument as in \cref{ch12:prop-corr}: $\rho>0$ inflates the fluctuations of $K$ by
$1+(N-1)\rho$ compared with the binomial. The distribution of $K$ also acquires a fat right tail and a fat mass at
$0$. Two standard ways to produce correlated Bernoulli indicators:

\begin{enumerate}[(a)]
  \item \emph{One-factor Gaussian model} (reconstruction; the model behind many credit-risk portfolios). Let
        $Z_i=\sqrt a\,Y+\sqrt{1-a}\,\varepsilon_i$ with $Y,\varepsilon_i$ i.i.d.\ $N(0,1)$ and
        declare project $i$ a success if $Z_i<c=\Phi^{-1}(p)$. Given the common factor $Y$ the projects are independent with
        success probability $q(Y)=\Phi\bigl((c-\sqrt a\,Y)/\sqrt{1-a}\bigr)$, so
        $\Prob(K\ge k)=\E_Y\bigl[\Prob(\mathrm{Bin}(N,q(Y))\ge k)\bigr]$, a one-dimensional integral. The Bernoulli
        correlation $\rho$ is smaller than the latent correlation $a$ (it is $0.026$ for $a=0.1$ and $0.058$ for $a=0.2$).
  \item \emph{Uncertainty about $p$} (beta-binomial). The lecturer's distinction between risk and uncertainty has an exact
        counterpart. Suppose the projects share an unknown success rate $\Pi\sim\mathrm{Beta}(\alpha,\beta)$ with
        $\E\Pi=p$ and, given $\Pi$, are independent Bernoulli($\Pi$). Then $\Cov(X_i,X_j)=\Var\Pi$ and
        \begin{equation}\label{ch12:eq-rhobeta}
          \rho=\frac{\Var\Pi}{p(1-p)}=\frac{1}{\alpha+\beta+1}.
        \end{equation}
        Indeed $\Var\Pi=p(1-p)/(\alpha+\beta+1)$ for a beta distribution. The parameter $\alpha+\beta$ is the
        ``effective sample size'' of the prior information about $p$: if the historical rate ($3.4\%$ against $5\%$)
        were known to only about 50 observations, $\alpha+\beta=49$ and $\rho=0.02$, a correlation that no amount of
        diversification can remove (\cref{ch12:tab-corr} with $\rho=0.02$).
\end{enumerate}

\begin{table}[ht]
\centering\small
\caption{Reconstruction: probability that the portfolio of $N$ projects with $p=5\%$ has at least $k$ successes,
independent case ($\rho=0$), one-factor Gaussian model with latent correlation $a$ and beta-binomial
($\alpha+\beta=49$, $\rho=0.02$). $N=150,k=3$ is break-even in the sense of the text.}\label{ch12:tab-counts}
\begin{tabular}{@{}lrrrr@{}}
\toprule
& $N=15,\ k\ge1$ & $N=150,\ k\ge1$ & $N=150,\ k\ge3$ & $N=150,\ k\ge4$\\
\midrule
independent          & $0.537$ & $0.9995$ & $0.982$ & $0.945$\\
Gaussian, $a=0.1$    & $0.482$ & $0.968$  & $0.821$ & $0.730$\\
Gaussian, $a=0.2$    & $0.432$ & $0.899$  & $0.689$ & $0.601$\\
Gaussian, $a=0.5$    & $0.300$ & $0.616$  & $0.419$ & $0.368$\\
beta-binomial        & $0.491$ & $0.970$  & $0.838$ & $0.753$\\
\bottomrule
\end{tabular}
\end{table}

\begin{intuition}[What correlation does to the tails]
The mean number of successes ($7.5$ out of 150) is unchanged; what changes is the probability of the unlucky
scenario in which \emph{everything fails together}. That scenario is invisible in the average and decisive for
someone who has to repay creditors. This is why the lecture insists on the near-independence of rare diseases
\ts{24}{18}{33:36} and why the portfolio manager's key parameter is a correlation, not a mean.
\end{intuition}

% ==================================================================
\section{From portfolio to structure: a securitisation sketch}\label{ch12:sec-tranche}

In the lecture the \$30 billion is obtained by appeal to the investing public \ts{24}{18}{30:30} and, for fusion, to
retail money \ts{24}{18}{57:45}. The natural instrument, which Lo calls a ``mega fund'' \ts{24}{18}{54:42} and
which the past final project of 2024 L1 develops for paediatric oncology, is a pool of projects funded by
tranches of different seniority. He does not describe tranches in this lecture; the following is a generic
\emph{Reconstruction}.

\begin{example}[Reconstruction: one senior debt tranche on the 150-project pool]\label{ch12:ex-tranche}
Pool of $N=150$ projects, all financed at year 0 by \$30 billion, with $K$ successes and terminal cash flow
$V=12.3\,K$ (billion dollars) at year 10. Suppose the whole \$30 billion is raised as zero-coupon debt at $5\%$,
so the face value due at year 10 is $30\times1.05^{10}=48.9$. It is repaid if and only if
$12.3K\ge48.9$, i.e.\ $K\ge4$. With independent projects $\Prob(K\ge4)=0.945$; in the one-factor model the same probability is
$0.73$ for $a=0.1$ and $0.60$ for $a=0.2$ (\cref{ch12:tab-counts}). A structure with a junior equity tranche that
absorbs the first losses, and only a fraction of the capital raised as senior debt, would push the
senior repayment probability far higher; but the equity holders take the correlation risk, and the
example shows that its price depends mostly on $\rho$, not on $N$.
\end{example}
The moral of \cref{ch12:ex-tranche} is that securitisation transforms one risk profile into several
(the same principle as in structured credit), and that a rating of the senior tranche is a statement about the joint
law of the project outcomes. The correlation input is exactly what the two rare-disease and fusion arguments in the
lecture are about.

% ==================================================================
\section{Where the data science and AI enter}\label{ch12:sec-ml}

Despite the title, the lecture says almost nothing about machine learning: artificial intelligence appears in
passing as a competitor for investors' money \ts{24}{18}{1:16:35}, and Lo's affiliation with CSAIL
opens the talk \ts{24}{18}{0:00}. His actual ``data science'' is a \emph{simulation}
of the portfolio with parameters supplied by industry experts \ts{24}{18}{35:38}. What follows is a
short \emph{Reconstruction} (not from the course) of how the data would enter a real model.

\begin{enumerate}[(a)]
  \item \emph{Success probability from features.} Model the outcome $Y_i\in\{0,1\}$ of project $i$ with covariates $x_i$
        (indication, mechanism, phase, sponsor, biomarker) as $\Prob(Y_i=1\mid x_i)=\sigma(w\T x_i)$, with
        $\sigma(u)=1/(1+e^{-u})$, and fit $w$ by minimising the log-loss
        $-\sum_i\bigl[Y_i\log\sigma(w\T x_i)+(1-Y_i)\log(1-\sigma(w\T x_i))\bigr]$ with a penalty on $\|w\|$
        (regularisation, \cref{ch:regularized}; more general classifiers, \cref{ch:ml}). The output is a
        project-specific $p_i$ in place of the pooled $p=5\%$; the portfolio mean is then $\sum_ip_iG_i$.
  \item \emph{Shrinkage and uncertainty.} With few historical trials, the estimate of $p$ is imprecise, and a Bayesian update
        of a $\mathrm{Beta}(\alpha,\beta)$ prior with $s$ successes in $n$ trials gives
        $\mathrm{Beta}(\alpha+s,\beta+n-s)$ with mean $(\alpha+s)/(\alpha+\beta+n)$. By \eqref{ch12:eq-rhobeta}
        the residual uncertainty about $p$ is the common factor that limits diversification; more data
        lowers $\rho$.
  \item \emph{Dependence between projects.} Correlations (shared mechanisms, shared regulators, shared
        macro conditions for funding) are estimated from the same features, e.g.\ by treating shared
        attributes as factors, as in \cref{ch:pca}; the sample correlation of a handful of binary outcomes
        is very noisy, so the structure of the factors matters more than the numbers.
\end{enumerate}

\begin{coursenote}[What the lecture actually contains]
No data set, code or slide of the lecture is available. The only quantitative statements are those reproduced in
\cref{ch12:sec-single,ch12:sec-portfolio}: $\$200$M, 10 years, $5\%$, \$12.3B, $12\%$, $423.5\%$, $150$, $34.6\%$, Sharpe
$\approx0.34$--$0.35$, the rare-disease reduction to 10--20 projects and \$400--500M, and a Sharpe ratio above 1 from a
Monte Carlo of a published paper. Everything else in this chapter is derived here.
\end{coursenote}

% ==================================================================
\section{Second example: a fusion mega fund}\label{ch12:sec-fusion}

The second half of the lecture applies the same idea to energy. The argument is that the ``energy transition''
is not a transition: consumption of oil, gas and coal is still growing as more people join the middle class,
so low-carbon, on-demand sources are needed in addition \ts{24}{18}{42:02}\ts{24}{18}{43:03}. The candidates listed are fusion, fission and
geothermal; fusion promises more energy per kilogram of fuel than fission and no long-lived waste, but requires the three
Lawson conditions of temperature, pressure (density) and confinement time \ts{24}{18}{47:17}\ts{24}{18}{48:20}. Lo points to
a 20-tesla magnet built at MIT with Commonwealth Fusion Systems in 2021 and a laboratory reaction with net energy
gain \ts{24}{18}{49:23}\ts{24}{18}{50:27}.

His question, asked with Dennis Whyte, is whether finance can pay for fusion. The idea is to \emph{decompose} ``a
commercially viable power plant'' into a portfolio of enabling technologies (sustained
plasma at about 100 million degrees, ultra-clean vacuum systems, magnetic confinement, and so on), each with a value of its
own, for instance a vacuum system also useful in drug development, and to treat the portfolio as in the biomedical
case, with perhaps 150 components \ts{24}{18}{51:34}\ts{24}{18}{52:37}\ts{24}{18}{53:39}. He adds a
board of scientific and business advisers whose participation would follow once serious capital is deployed
\ts{24}{18}{54:42}.

A student asks whether the components would not be highly correlated. The answer is twofold: complete reactors from
different companies are correlated, but not perfectly, because they use different approaches (magnetic versus
inertial); and, more importantly, technologies such as heating, vacuum and magnet shaping have little to do with each
other \ts{24}{18}{56:44}\ts{24}{18}{57:45}. In the language of \cref{ch12:prop-corr}: the argument is that
$\rho$ is small, and it should be treated as the key empirical input, since the floor $\mathrm{SR}_1/\sqrt\rho$
is not eliminated by adding components that share a common ``fusion is feasible'' factor (compare \eqref{ch12:eq-rhobeta}).

\begin{coursenote}[Erratum: the retail-funding arithmetic]
Lo asks for a one-time \$1,500 from each of $3\%$ of 130 million households, ``more than \$30 billion''
\ts{24}{18}{56:44}\ts{24}{18}{57:45}. The product is
\[
  130\times10^{6}\times0.03\times1500=\$5.85\times10^{9},
\]
about \$5.9 billion, not \$30 billion; reaching \$30 billion at \$1,500 each needs about $15\%$ of households.
(He also says once that fusion is ``when we split an atom'': that is fission; fusion joins nuclei. Both are slips of
speech in the transcript; the intended meaning is clear.)
\end{coursenote}

\section{The rest of the lecture}\label{ch12:sec-rest}

\paragraph{The Harvey Lodish story} \ts{24}{18}{58:45}. As a young faculty member in 1983 Lodish was
approached to help develop a treatment for Gaucher's disease, a genetic enzyme deficiency; the company became Genzyme,
which got the first enzyme-replacement therapy approved in 1991 and was sold to Sanofi for \$20 billion in 2014.
In 2002 his unborn grandson was diagnosed with the same disease, and the drug saved him. The point Lo draws is that finance need not be
zero-sum: ``doing well by doing good'' \ts{24}{18}{1:02:57}.

\paragraph{Questions.}
\begin{itemize}
  \item \emph{Patents} \ts{24}{18}{1:05:05}: a patent gives exclusive use of the compound for 20 years (which is why revenues
        run over years 11--20 in \cref{ch12:sec-single}); afterwards anyone can manufacture cheaply, like aspirin.
        Producing the pill is cheap; the expensive part is the trials that lead to approval.
  \item \emph{How to find and evaluate opportunities} \ts{24}{18}{1:07:08}\ts{24}{18}{1:12:24}: expertise. Neil Kumar (BridgeBio's
        co-founder, a chemical-engineering PhD) worked at a consultancy on pharma strategy, then at a biotech venture
        firm, which had passed on single-drug projects as too risky; the simulations gave him the conviction to fund
        them himself.
  \item \emph{How can students invest?} \ts{24}{18}{1:08:11}: BridgeBio is public, at very high risk; private companies are
        harder; MIT is rich in opportunities for investors; Lo also teaches health-care financing (15.482, with an online
        version in MIT's Open Learning Library) \ts{24}{18}{1:11:21}.
  \item \emph{Crises} \ts{24}{18}{1:15:32}: COVID-19 temporarily removed the valley of death: the time between the
        sequencing of the virus and the first human injection was 63 days, against 3--5 years normally; afterwards the valley returned
        because vaccines are not attractive to investors compared with, for example, AI \ts{24}{18}{1:17:39}\ts{24}{18}{1:18:40}.
\end{itemize}

% ==================================================================
\begin{keyformulas}
\begin{itemize}
  \item Sharpe ratio: $\mathrm{SR}=(\E R-r_f)/\sqrt{\Var R}$.
  \item Single project (lecture): cost \$200M, 10 years, $p=5\%$, payoff \$12.3B; $\E G=3.075$ per dollar,
        $\sqrt{\Var G}=61.5\sqrt{p(1-p)}=13.4$ over 10 years ($424\%$ per year after dividing by $\sqrt{10}$).
  \item $N$ i.i.d.\ projects: $\Var R_P=\sigma^2/N$, $\mathrm{SR}_N=\sqrt N\,\mathrm{SR}_1$.
  \item Common correlation $\rho$: $\Var R_P=\sigma^2\bigl(\tfrac1N+\tfrac{N-1}{N}\rho\bigr)$,
        $\mathrm{SR}_N\to\mathrm{SR}_1/\sqrt\rho$.
  \item Successes: $\Var K=Np(1-p)\bigl(1+(N-1)\rho\bigr)$; $\Prob(K\ge1)=1-(1-p)^N$ if independent.
  \item Unknown $p\sim\mathrm{Beta}(\alpha,\beta)$: $\rho=1/(\alpha+\beta+1)$.
\end{itemize}
\end{keyformulas}

\section*{Exercises}
\addcontentsline{toc}{section}{Exercises}

\begin{exercise}[not from the course]
Verify \eqref{ch12:eq-annuity}: find the discount rate for which \$2 billion in each of years 11--20 is worth
\$12.3 billion at year 10, and compute the value at year 0 (in units of the \$200M cost) of the project of
\cref{ch12:ex-single}, using its expected payoff.
\end{exercise}

\begin{exercise}[not from the course]
Prove \eqref{ch12:eq-varK}. Then compute the standard deviation of $K$ for $N=150$, $p=0.05$ and $\rho\in\{0,0.02,0.05\}$
and compare it with the mean $7.5$.
\end{exercise}

\begin{exercise}[not from the course]
With $R_P=\frac1N\sum_iR_i$ and general covariance matrix $\Sigma$, the variance is $\frac{1}{N^2}\mathbf 1\T\Sigma\mathbf 1$
(\cref{ch:portfolio}). Show that for the common-correlation matrix $\Sigma=\sigma^2[(1-\rho)I+\rho\,\mathbf 1\mathbf 1\T]$ this reproduces
\eqref{ch12:eq-varcorr}, and find its eigenvalues. Which eigenvector carries the ``undiversifiable'' part?
\end{exercise}

\begin{exercise}[not from the course]
Show that in the beta-binomial model of \cref{ch12:sec-counts}(b), $\rho=1/(\alpha+\beta+1)$, and that after observing
$s$ successes in $n$ historical trials of a common process the posterior correlation between two future projects is
$1/(\alpha+\beta+n+1)$. What sample size $n$ makes it fall below $0.005$ starting from $\alpha+\beta=49$?
\end{exercise}

\begin{exercise}[not from the course]
In \cref{ch12:ex-tranche} let the senior debt be only a fraction $f$ of the \$30 billion, with face value $30f\times1.05^{10}$.
Find, for the independent model, the smallest $f$ such that the senior tranche is repaid with probability at least $0.99$.
(The rest of the capital is equity and absorbs the first losses.)
\end{exercise}

\begin{exercise}[not from the course]
Erratum check: with the numbers of \cref{ch12:sec-fusion}, what fraction of the 130 million households must contribute \$1,500
to raise \$30 billion? What one-time contribution per household would be needed from $3\%$ of them?
\end{exercise}

\begin{exercise}[not from the course]
Simulate the one-factor Gaussian model of \cref{ch12:sec-counts}(a) for $N=150$, $p=0.05$, $a=0.1$ and estimate
$\Prob(K\ge3)$ by Monte Carlo with $10^5$ paths; compare with \cref{ch12:tab-counts}.
\end{exercise}
